% Generated draft: requires mathematical and source review.
\documentclass[11pt]{article}
\usepackage[margin=25mm]{geometry}
\usepackage{fontspec}
\setmainfont{Latin Modern Roman}
\usepackage{amsmath,amssymb,mathtools,unicode-math}
\setmathfont{Latin Modern Math}
\usepackage{enumitem,needspace,xurl,hyperref}
\setlength{\parindent}{0pt}
\setlength{\parskip}{5pt}
\setlength{\emergencystretch}{3em}
\newcommand{\N}{\mathbb{N}}
\newcommand{\Z}{\mathbb{Z}}
\newcommand{\Q}{\mathbb{Q}}
\newcommand{\R}{\mathbb{R}}
\newcommand{\C}{\mathbb{C}}
\newcommand{\F}{\mathbb{F}}
\newcommand{\A}{\mathbb{A}}
\newcommand{\PP}{\mathbb{P}}
\newcommand{\E}{\mathbb{E}}
\newcommand{\eps}{\varepsilon}
\newcommand{\dd}{\,\mathrm{d}}
\begin{document}
\section*{UnsolvedMath ID: 701341}
\textbf{Source problem code:} catalogue entry
\subsection{Definitions and mathematical statement}
Let \(G\) be a finite soluble group, meaning that its derived series terminates at the identity subgroup. Write \(\pi(G)\) for the set of primes dividing \(|G|\). A Sylow complement basis of \(G\) is a family \(\mathcal B=(H_p)_{p\in\pi(G)}\) of pairwise permuting subgroups, where \(H_p\) is a Hall \(p'\)-subgroup: its order is relatively prime to \(p\), and its index in \(G\) is a power of \(p\). The associated \emph{system normalizer} is
\[
D(\mathcal B)=\bigcap_{p\in\pi(G)}N_G(H_p),
\qquad N_G(H)=\{g\in G:gHg^{-1}=H\}.
\]
Let \(\mathcal D(G)\) denote the set of system normalizers obtained from all Sylow complement bases of \(G\). For the trivial group, the intersection over the empty prime set is understood to be \(G\).

Define the upper central series by
\[
Z_0(G)=1,\qquad Z_{i+1}(G)/Z_i(G)=Z\bigl(G/Z_i(G)\bigr)\quad(i\geq 0).
\]
The \emph{hypercenter} is \(Z_\infty(G)=\bigcup_{i\geq0}Z_i(G)\); this series stabilizes because \(G\) is finite.

The question attributed to A. F. Vasil'ev and L. A. Shemetkov is whether
\[
\exists n\in\Z_{>0}\ \forall\text{ finite soluble groups }G\
\exists D_1,\ldots,D_n\in\mathcal D(G):\qquad
Z_\infty(G)=\bigcap_{j=1}^{n}D_j.
\]
Thus \(n\) must be independent of \(G\), while the chosen system normalizers may depend on \(G\). If such an integer exists, determine the least one. The displayed formulation permits repeated system normalizers, equivalently asking for an intersection of at most \(n\) system normalizers. The source does not explicitly address distinctness, so this convention requires review.

The source's \emph{Comment of 2017} reports that the assertion is proved for groups with trivial Frattini subgroup, where \(\Phi(G)\), the intersection of the maximal proper subgroups of \(G\), is trivial. It cites a 2018 journal article but supplies neither a numerical bound nor an explicit description of which part of the question is settled there. The supplied registry labels the general problem open and separately records an explicitly unverified AI claim that no such integer exists. That claim is not adopted as a mathematical conclusion.
\subsection{Short English statement}
Is there a single positive integer \(n\) such that the hypercenter of every finite soluble group is the intersection of \(n\) suitably chosen system normalizers? If so, find the least \(n\). Repetitions are permitted in this draft, subject to review of the source wording. The source reports a result for groups with trivial Frattini subgroup and an explicitly unverified negative AI claim for the general question.
\subsection{Sources}
\begin{itemize}
\item Local UnsolvedMath registry, supplied complete record with permanent ID \(701341\), exact source code \texttt{catalogue entry}, title \emph{Kourovka Notebook Problem 17.39}. Proposed by A. F. Vasil'ev and L. A. Shemetkov in 2010. Upstream revision: \texttt{d371ee917ac9647fbb06cd6fc529deec6f3dda82}.
\item \emph{The Kourovka Notebook (Unsolved Problems in Group Theory)}, original problem 17.39, 17th issue (2010); also identified in the supplied record with the 21st issue edition (2026), \texttt{arXiv:1401.0300v48}. The supplied collection description names E. I. Khukhro and V. D. Mazurov as editors.
\item The supplied \emph{Comment of 2017} cites A. Ballester-Bolinches, J. Cossey, S. F. Kamornikov, and H. Meng, \emph{J. Algebra} \textbf{499} (2018), 438--449, for the reported result concerning groups with trivial Frattini subgroup. No article title is supplied.
\item The supplied background and status review attribute an explicitly unverified AI solution claim to M. Kinyon and J. Urban, \emph{Preprint} (2026): \url{https://kourovkanotebookorg.wordpress.com/wp-content/uploads/2026/09/kourovka-mathematics.pdf}. The reported claim is that no such integer exists; it is not treated here as an established result.
\item The supplied status review is dated \texttt{2026-10-16}, labels the canonical status \texttt{open}, and classifies its evidence as \texttt{unrefereed\_claim}. It reports a \# marker for the unverified AI claim in the October 2026 21st issue, \texttt{arXiv:1401.0300v48}. This is supplied dated curation, not an independently verified literature assessment. None of the cited sources has been independently consulted for this draft.
\end{itemize}
\end{document}
